%HOMEWORK template for M 325K
\documentclass{article}
\headheight=7pt         \topmargin=14pt
\textheight=574pt       \textwidth=445pt
\oddsidemargin=18pt     \evensidemargin=18pt

%Begin package import

\usepackage{amsmath, amssymb, amsthm}
\usepackage{mathtools}
\usepackage{pdfpages}
\usepackage{tikz-cd}

%End package import
%Begin custom commands

\newcommand{\C}{\mathbb{C}}
\newcommand{\R}{\mathbb{R}}
\newcommand{\Q}{\mathbb{Q}}
\newcommand{\Z}{\mathbb{Z}}
\newcommand{\N}{\mathbb{N}}
\newcommand{\F}{\mathbb{F}}
\newcommand{\abs}[1]{\left\lvert #1\right\rvert}
\newcommand{\RM}[1]{\mathrm{#1}}
\newcommand{\BF}[1]{\textbf{#1}}
\newcommand{\e}{\epsilon}
\newcommand{\ceil}[1]{\lceil{#1}\rceil}
\newcommand{\floor}[1]{\lfloor{#1}\rfloor}
\newcommand{\rarr}[0]{\rightarrow}
\newcommand{\IM}[1]{\text{Im}(#1)}
\newcommand{\Hom}[0]{\text{Hom}}
\newcommand{\Pow}[1]{\text{Pow}(#1)}
\newcommand{\angles}[1]{\langle #1 \rangle}

%End custom commands
%Begin custom theorems

\newtheorem{innercustomgeneric}{\customgenericname}
\providecommand{\customgenericname}{}
\newcommand{\newcustomtheorem}[2]{%
\newenvironment{#1}[1]
{%
\renewcommand\customgenericname{#2}%
\renewcommand\theinnercustomgeneric{##1}%
\innercustomgeneric%
}
{\endinnercustomgeneric}
}

\newcustomtheorem{THRM}{Theorem}
\newcustomtheorem{LEM}{Lemma}
\newcustomtheorem{EX}{Exercise}
\newcustomtheorem{COR}{Corollary}
\newcustomtheorem{PROB}{Problem}
\newcustomtheorem{claim}{Claim}

\newenvironment{solution}
{\renewcommand\qedsymbol{$\blacksquare$}\begin{proof}[Solution]}
{\end{proof}}

\newenvironment{definition}
{\renewcommand\qedsymbol{$\blacksquare$}\begin{proof}[Definition]}
{\end{proof}}

%End custom theorems
\begin{document}

\title{Homework 12}
\author{Arham Lodha}%put your name here
\date{April 16, 2024}%enter the due date here
\maketitle

\begin{PROB}{1a}\label{Problem 1a.}
\end{PROB}
\begin{proof}
    $\forall \epsilon > 0$. Let $\delta = \epsilon / 2$. $\forall x, u \in [1, \infty) \ni \abs{x - u} < \delta$.  
    \begin{align*}
        \abs{ \frac{1}{x^2} - \frac{1}{u^2} } &= \abs{\frac{u^2 - x^2}{x^2 u^2}}
        \\&= \abs{\frac{(u - x)(u + x)}{x^2u^2}}
        \\&= \delta \abs{\frac{u + x}{x^2u^2}}
        \\&= \delta \abs{\frac{1}{x^2 u} + \frac{1}{x u^2}}
        \\&< 2\delta \text{ (since $x, u \geq 1$)}
        \\&= \epsilon
    \end{align*}

    Thus $\frac{1}{x^2}$ is uniformly continous on $[1, \infty)$  
\end{proof}
\begin{PROB}{}\label{Problem .}
\end{PROB}
\begin{proof}
    Let $\epsilon_0 = 0.5$.  
    Let $(x_n) = \frac{1}{\sqrt{n}}$ and let $(u_n) = \frac{1}{\sqrt{n + 1}}$. $x_n$ and $u_n$ converge to 0. By operation of limits, $\lim_{n \to \infty}(x_n - u_n) = 0$. However $\forall n \in \N$, 
    
    \begin{align*}
        \abs{\frac{1}{x_n^2} - \frac{1}{u_n^2}} &= \abs{\frac{1}{\frac{1}{(\sqrt{n})^2}} - \frac{1}{\frac{1}{(\sqrt{n + 1})^2}}} \\
        &= \abs{n - n - 1} \\
        &= 1 > \frac{1}{2}
    \end{align*}

    Thus by Nonuniform convergence criteron, $\frac{1}{x^2}$ has uniform convergence on $(0, \infty)$.  
\end{proof}

\bigskip

\begin{PROB}{2}\label{Problem 2.}
\end{PROB}
\begin{proof}
    By the definition of a closed set, $I = \overline{I} = I \cup I'$ and $I'$ is the set of limit points of I. Thus I has the set of limit points of $I$. Thus since $x_0$ is the limit of a convergent sequence $(x_n) \subset E\subset I$, $x_0 \in I$. 
    
    By the definition of sequential definition of continuity, since $x_n \to x_0$, $f(x_n) \to f(x_0)$, and $g(x_n) \to g(x_0)$. Thus, $\lim_{n \to \infty} f(x_n) = f(x_0)$ and $\lim_{n \to \infty} g(x_n) = g(x_0)$ , however since $\forall n \in \N$, $x_n \in E$ thus $f(x_n) = g(x_n)$. Thus $\lim_{n \to \infty} g(x_n) = f(x_0)$. By the uniqueness of the limit, $f(x_0) = g(x_0)$. Thus $x_0 \in E$.              
\end{proof}

\bigskip

\begin{PROB}{3}\label{Problem 3.}
\end{PROB}
\begin{proof}
    Let $(x_n)$ be the following recursive sequence. $x_1 = a$. $\forall n \in \N$, $\exists y \in I \ni \abs{f(y)} \leq \frac{1}{2}\abs{f(x_n)}$, let $x_{n + 1} = y$. Thus by construction, $\forall n \in \N$, $\abs{f(x_{n + 1})} \leq \frac{1}{2}\abs{f(x_n)} \leq \dots \leq \frac{1}{2^{n - 1}}\abs{f(x_1)} = \frac{1}{2^{n - 1}} \abs{f(a)}$.   

    $(x_n) \subset I \implies \forall n \in \N, a \leq x_n \leq b$. Thus $x_n$ is bounded. By Bolzano Wierstrass theorem, $\exists y_n$ a subsequence of $x_n$ which is monotonic, since $x_n$ is bounded. Since $y_n$ is a subsequence of $x_n$ it is bounded as well, by monotonic subsequence theorem, $y_n$ is a convergent sequence. Thus $\exists c \in \R \ni y_n \to c$. By the definition of a closed set, $I = \overline{I} = I \cup I'$ and $I'$ is the set of limit points of I. Thus I has the set of limit points of $I$. Thus since $c$ is the limit of a convergent sequence $(y_n) \subset E\subset I$, $c \in I$.
    
    Since $\forall n \in \N$, $y_n \in (x_n)$ and by construction, $\forall n \in \N$, $\abs{f(y_{n + 1})} \leq \frac{1}{2}\abs{f(y_n)} \leq \dots \leq \frac{1}{2^{n - 1}}\abs{f(y_1)} \leq \frac{1}{2^{n - 1}} \abs{f(a)}$.
    
    Thus, $0 \leq \abs{f(y_n)} \leq \frac{1}{2^{n - 1}}\abs{f(a)}$. 
    
    $\forall \epsilon > 0$. By the Archmedian property, $\exists K \in \N \ni K > -\log_2(\epsilon)$. $\forall x \geq K$.   
    
    \begin{align*}
        \abs{\frac{1}{2^n}} &= \frac{1}{2^n} \\
        &= \frac{1}{2^K} \text{(since $n \geq K \implies 2^n \geq 2^K$) }\\
        &< \epsilon 
    \end{align*}

    Thus $\lim_{n \to \infty} \frac{1}{2^n} = 0$. By operation of limits, $\lim_{n \to \infty} \frac{2}{2^{n}} \abs{f(a)} = \lim_{n \to \infty} \frac{1}{2^{n - 1}} \abs{f(a)} = 0$. Limit of constant sequence is constant, thus 0 converges to 0. Thus by squeeze theorem, $\lim_{n \to \infty} f(y_n) = 0$. By the sequential definition of continuity, $\lim_{n \to \infty} f(y_n) = f(c)$. By the uniquess of the limit, $f(c) = 0$. Thus $\exists c \in I \ni f(c) = 0.$      
\end{proof}

\bigskip

\begin{PROB}{4}\label{Problem 4.}
\end{PROB}
\begin{proof}
    Let $f(x) = x$, $g(x) = c$ .  
    $\forall r \in \R$. $\forall \epsilon > 0$, let $\delta = \epsilon$. $\forall x \in \R \ni \abs{x - r} < \delta$. Thus $\abs{f(x) - f(r)} = \abs{x - r} < \delta = \epsilon$. Let $\gamma = 1$. $\forall x \in \R \ni \abs{x - r} < \gamma$. $\abs{g(x) - g(r)} = \abs{c - c} = 0 < \epsilon$.     Thus, $f(x), g(x)$ are continous at $\forall r \in \R$. Thus $f(x), g(x)$ are continous on $\R$. $\forall n \in \N$, $a_n x^n = a_n(\prod_{k = 0}^{n}f)(x)$, thus by operation of limits, $a_n x^n$ is continous. Thus $p(x) = \sum_{n = 0}^{\deg p}a_n x^n$ is continous by operation of limits. 
    
    Thus $p(x) = x^4 + 7x^3 - 9$ is continous. $p(0) = -9$, $p(2) = 63$, $p(-10) = 10^4 - 7(10^3) - 9 = 10000 - 7000 - 9 = 3000 - 9 = 2997$ .  By the intermediate value theorem, $\exists a \in (0, 2) \ni p(c) = 0$ and $\exists b \in (-10, 0) \ni p(b) = 0$. Since $(-10, 0) \cap (0, 2) = \emptyset$, $a,b$ are distinct. Thus p has at least 2 roots. 

    (Note a alternate proof would use FT of Algebra and the fact complex roots come with its conjugates, so it suffices to show there is at least 1 real root.)
\end{proof}

\end{document}